% Chapter 8: CLAUDE.md Architecture
\chapter{CLAUDE.md Architecture}
\label{ch:architecture}

\begin{itshape}
You have five projects. Each has its own CLAUDE.md with 80\% identical content.
One project's CLAUDE.md has grown to 600 lines and Claude ignores half the rules.
You need architecture, not just configuration.
\end{itshape}

In \cref{ch:first-claude-md}, you created your first CLAUDE.md.
Now that you have experience with multiple projects, we will explore
the architecture behind configurations that scale.

\section{The Hub-and-Spoke Pattern}
\label{sec:hub-and-spoke}

\practitioner{For developers managing multiple repositories,
maintain a central ``hub'' with shared patterns referenced from each project.}

\begin{fullwidth}
\begin{beforeafter}
\before
Five CLAUDE.md files, each 200 lines, 80\% identical.
Updating git conventions means editing five files.
Three of them are out of date.

\after
One hub with shared patterns. Five spokes, each 40 lines.
Update once, all projects inherit.

\begin{minted}{markdown}
# Project CLAUDE.md (spoke — 40 lines)
@~/shared-patterns/git.md
@~/shared-patterns/testing.md

## Project-Specific
- Build: `make digital`
- Test: `pytest tests/`
- Architecture: Tufte-style LaTeX, one chapter per file
\end{minted}
\end{beforeafter}
\end{fullwidth}

\official{Use the \code{@path/to/file.md} import syntax to reference
external files without duplicating content.\sourceurl{https://code.claude.com/docs/en/memory}}

\begin{whybox}[Why Hub-and-Spoke]
Duplication across CLAUDE.md files is not just inconvenient---it is
actively harmful. When the git convention changes in one file but not others,
Claude learns different conventions depending on which project it is in.
Inconsistency accumulates until you cannot trust that any project's
instructions are current.
\end{whybox}

% -------------------------------------------------------------------
\section{Tier Classification}
\label{sec:tiers}

Not every project needs the same level of configuration.
Classify by complexity:

\begin{fullwidth}
\begin{tabular}{@{}llp{8cm}@{}}
\toprule
\textbf{Tier} & \textbf{Components} & \textbf{When to Use} \\
\midrule
T0 & None & Prototypes, learning repos \\
T1 & CLAUDE.md only & Small scripts, documentation projects \\
T2 & CLAUDE.md + settings & Standard development projects \\
T3 & + skills, commands, rules & Multi-module projects with workflows \\
T4 & + hooks, MCP & Projects requiring automation and quality gates \\
T5 & + agent teams, pipelines & Enterprise-scale projects with CI/CD \\
\bottomrule
\end{tabular}
\end{fullwidth}

\margintip{Start at T1. Promote when friction justifies it.
Most solo projects stabilize at T2--T3.}

% -------------------------------------------------------------------
\section{Conditional Rules}
\label{sec:conditional-rules}

\official{Place \code{.md} files in \code{.claude/rules/} with \code{paths:}
frontmatter to control when they load.\sourceurl{https://code.claude.com/docs/en/best-practices}}

Rules without \code{paths:} frontmatter load on \emph{every} session.
For monorepos and multi-module projects, this bloats context with
irrelevant instructions.

\begin{minted}{markdown}
---
paths: notebooks/**/*.py
---
# Notebook Rules
- No hardcoded file paths (use config)
- Every notebook has a markdown header explaining purpose
- Run `nbstripout` before committing
\end{minted}

\begin{minted}{markdown}
---
paths: src/models/**/*.py
---
# Model Training Rules
- Log all hyperparameters to MLflow
- Save model artifacts to models/
- Include reproducibility seed in all configs
- Run `pytest tests/models/` after changes
\end{minted}

\begin{warnbox}[Common Mistake]
Forgetting \code{paths:} frontmatter. The rule loads globally,
Claude's context fills with irrelevant instructions, and important
rules from CLAUDE.md get deprioritized.
\end{warnbox}

% -------------------------------------------------------------------
\section{Settings: 5-Level Precedence}
\label{sec:settings}

\marginnote[Official]{Settings documentation:\sourceurl{https://code.claude.com/docs/en/settings}}

Claude Code uses a 5-level settings hierarchy.
Higher levels override lower:

\begin{enumerate}
  \item \textbf{Managed} (\code{managed-settings.json}) --- IT-deployed, highest priority
  \item \textbf{CLI arguments} --- temporary session overrides
  \item \textbf{Local} (\code{.claude/settings.local.json}) --- personal, auto-gitignored
  \item \textbf{Project} (\code{.claude/settings.json}) --- team standards, committed to git
  \item \textbf{User} (\code{\textasciitilde/.claude/settings.json}) --- global personal defaults
\end{enumerate}

\practitioner{Convention: commit \code{settings.json} for team-wide standards.
Keep \code{settings.local.json} for personal preferences (model choice, output style).}

% -------------------------------------------------------------------
\section{CLAUDE.md Length and Attention}
\label{sec:length}

\marginwarning{Keep core CLAUDE.md under $\sim$300 lines. Beyond that, attention per rule
tends to decrease. This is heuristic, not a hard threshold---but the effect is real.
(Total config including \code{@imports} may reach 500 lines safely since imports load
only when relevant.)}

\practitioner{When CLAUDE.md grows long:}

\begin{enumerate}
  \item Move module-specific rules to \code{.claude/rules/} with \code{paths:}
  \item Move shared patterns to a hub via \code{@path} imports
  \item Move enforcement rules to hooks (they cannot be ignored)
  \item Keep core CLAUDE.md to project-wide essentials: build commands,
    architecture, verification criteria
\end{enumerate}

% -------------------------------------------------------------------
\section{Output Styles}
\label{sec:output-styles}

\marginnote[Official]{Configure via \code{outputStyle} in settings
or \code{--output-style} flag.}

\official{Output styles customize how Claude structures its responses:
tone, verbosity, formatting conventions, and problem-solving protocols.}

\practitioner{Custom styles are most valuable when they encode
\emph{decision protocols}, not just formatting preferences.}
A style that says ``use bullets'' saves minimal effort.
A style that says ``classify the issue type before searching for solutions''
changes the quality of every interaction.

% -------------------------------------------------------------------
\section{Visual: The Configuration Hierarchy}

\begin{figure}[H]
\begin{fullwidth}
\centering
\begin{tikzpicture}[
  layer/.style={draw=WarmBlue, fill=WarmBlue!8, rounded corners=3pt,
    minimum width=10cm, minimum height=1.1cm, font=\small, align=center},
  label/.style={font=\footnotesize\itshape, text=WarmBlue!70!black},
  >=Stealth
]
  \node[layer] (output) at (0,0) {\textbf{Output Styles} --- Communication patterns, tone, protocols};
  \node[layer] (settings) at (0,1.5) {\textbf{Settings} --- Permissions, model, hooks, env vars};
  \node[layer] (rules) at (0,3.0) {\textbf{Rules} --- Conditional context (paths: frontmatter)};
  \node[layer] (claude) at (0,4.5) {\textbf{CLAUDE.md} --- Project knowledge, build commands, architecture};

  \draw[->, thick, WarmBlue!60] (5.5,0) -- (5.5,4.5)
    node[midway, right, label] {Higher priority};

  \node[label, anchor=east] at (-5.3,4.5) {Global + Project};
  \node[label, anchor=east] at (-5.3,3.0) {Per-module};
  \node[label, anchor=east] at (-5.3,1.5) {5-level cascade};
  \node[label, anchor=east] at (-5.3,0) {Per-user};
\end{tikzpicture}
\caption{The four configuration layers, each handling a distinct concern.}
\label{fig:config-hierarchy}
\end{fullwidth}
\end{figure}

\begin{keyinsight}
The configuration hierarchy separates concerns:
CLAUDE.md for project knowledge, rules for conditional context,
settings for permissions and behavior, output styles for communication.
Mixing these concerns (e.g., putting permission rules in CLAUDE.md)
reduces their effectiveness.
\end{keyinsight}

% -------------------------------------------------------------------
\begin{trybox}
If you have 3+ projects:
\begin{enumerate}
  \item Create a shared patterns directory (\code{\textasciitilde/shared-patterns/}).
  \item Move your git commit conventions and testing standards there.
  \item Replace the duplicated content in each project's CLAUDE.md
    with \code{@\textasciitilde/shared-patterns/git.md} imports.
  \item Verify: each project's CLAUDE.md should be under 100 lines.
\end{enumerate}
If you have one project: count the lines in your CLAUDE.md.
If over 300, identify 3 sections that should be conditional rules.
\end{trybox}
